\documentclass[11pt]{ctexart}

\usepackage{geometry}
\geometry{a4paper, margin=2.5cm}

\usepackage{hyperref}
\usepackage{enumitem}
\usepackage{longtable}
\usepackage{array}

\title{Snooker2D 中期分报告（开发者 B）}
\author{Ironhammer Std}
\date{2026年7月13日}

\begin{document}

\maketitle

\section{负责范围}

Model 层（游戏状态与规则引擎）和 ViewModel 层（MVVM 桥接）。负责文件：

\begin{center}
\begin{tabular}{|l|l|}
\hline
\textbf{文件} & \textbf{说明} \\
\hline
\texttt{src/model/Ball.h/.cpp} & 球的属性：类型、位置、速度、是否落袋 \\
\texttt{src/model/Table.h/.cpp} & 球桌：尺寸、库边定义、袋口坐标 \\
\texttt{src/model/Physics.h/.cpp} & 物理引擎：碰撞检测、摩擦力、袋口检测 \\
\texttt{src/model/GameState.h/.cpp} & 核心状态机：开球、击球、回合切换 \\
\texttt{src/model/Player.h/.cpp} & 玩家：姓名、得分 \\
\texttt{src/model/Rules.h/.cpp} & 规则模块：犯规检测、罚分计算 \\
\texttt{src/viewmodel/GameViewModel.h/.cpp} & 核心 ViewModel：属性暴露、命令实现 \\
\texttt{src/viewmodel/CueControlViewModel.h/.cpp} & 击球控制：角度、力度、加塞 \\
\texttt{src/viewmodel/ScoreViewModel.h/.cpp} & 计分板数据：双方累计分、单杆分 \\
\hline
\end{tabular}
\end{center}

\section{已完成工作}

\subsection{Model 层数据结构（7/10）}

\begin{itemize}
\item \textbf{Ball}：位置（Vector2D）、速度、类型（BallType 枚举）、落袋状态、台面状态。22 颗球按斯诺克标准摆法初始化
\item \textbf{Table}：800×400 台面、6 条库边（Cushion 结构体）、6 个袋口（Pocket 结构体）
\item \textbf{Player}：得分属性（scoreChanged 信号）、姓名
\end{itemize}

\subsection{物理引擎（7/12）}

\begin{enumerate}[leftmargin=*]
\item 匀速直线运动（moveBalls：位置 += 速度 × deltaTime）
\item 摩擦力减速（applyFriction：速度 *= FRICTION\_COEFFICIENT，0.99/帧）
\item 速度阈值停止（MIN\_VELOCITY = 0.05，低于此值设为零）
\item 球-球弹性碰撞（resolveBallCollision）：动量守恒碰撞公式
\item 球-库边反弹（resolveCushionCollision）：通过 closestPointOnSegment 计算库边接触点，按库边弹性系数反射
\item 袋口检测（checkPocketDetection）：球心与袋口距离 < POCKET\_RADIUS 判定落袋
\end{enumerate}

物理参数：碰撞恢复系数 0.88，库边恢复系数 0.77，摩擦系数 0.99，固定时间步长 1/60 秒。

\subsection{规则引擎初版（7/12--7/13）}

\begin{itemize}
\item FoulResult 结构体：犯规类型 + 描述文本 + 罚分
\item 犯规检测框架：空杆、打错球、白球落袋、无球碰库
\item 罚分计算：取被犯规球分值，最少罚 4 分
\item 白球落袋后复位逻辑
\item 开局黑球位置重叠修复
\item 分数处理与回合转移
\end{itemize}

\subsection{ViewModel 桥接}

\begin{itemize}
\item GameViewModel：Q\_PROPERTY 暴露球位置列表、当前玩家、游戏阶段、比分、角度/力度；Q\_INVOKABLE 命令 shoot/setAngle/setPower/confirmFoul/selectFreeBall
\item CueControlViewModel：角度/力度属性、shootRequested 信号、reset/setAngleAndPower 方法
\item ScoreViewModel：双方得分、单杆分、犯规消息、状态消息属性
\item Model 信号 → ViewModel 属性同步：onModelPhaseChanged / onModelTurnChanged / onModelSimulationFinished / onModelFoulOccurred
\end{itemize}

\subsection{跨层协调}

\begin{itemize}
\item 在 MainWindow 中补充角度/力度实时同步连接（CueControlVM ↔ GameVM）
\item 配合 A 完成击球流程从按钮到鼠标的重构
\item 新增 GameViewModel::gameRestarted 信号，支持比赛重启时的状态清理
\end{itemize}

\section{提交统计}

共 5 条提交，日期 2026-07-12 至 2026-07-13。工作集中在中期检查前两日。

\section{技术要点}

\begin{itemize}
\item 物理引擎采用固定时间步长（1/60s），避免帧率波动影响模拟精度
\item 球-球碰撞使用弹性碰撞公式，基于两球心连线方向交换速度分量
\item 库边碰撞通过 closestPointOnSegment 计算接触点法线，精度优于简单 AABB 检测
\item Model 层所有类继承 QObject，使用 Qt 信号/槽通知状态变更
\item ViewModel 不直接耦合 View 层，通过 Q\_PROPERTY NOTIFY 信号推送更新
\end{itemize}

\section{遇到的技术难点}

\subsection{物理参数的经验性调优}

碰撞恢复系数、库边反弹系数和摩擦系数无法通过理论计算直接确定——真实斯诺克的物理参数受台呢材质、球弹性、环境温度等因素影响。初始参数（恢复系数 0.96、摩擦 0.98）导致球碰撞后反弹过度，贴库球几乎无法静止。通过逐步降低恢复系数（0.96→0.90→0.88）、增加库边独立的弹性系数（0.77，显著低于球-球），以及提高摩擦（0.98→0.99），最终达到可接受的运动效果。这一过程消耗了大量试错时间，也说明物理模拟的参数工程与算法本身同样重要。

\subsection{库边碰撞检测的精度问题}

简单的 AABB 矩形检测无法正确处理球撞击库边角点的情况——球可能在库边的端点附近出现穿透或抖动。解决方案是引入 closestPointOnSegment 算法：计算球心到库边线段的最近点，若距离小于球半径则判定为碰撞，法线方向为球心指向最近点。该算法由 C 在 Common 层实现，B 在 Physics 层调用。

\subsection{ViewModel 属性的同步时机}

GameViewModel 需要在多个信号回调中刷新球位置和分数（onModelPhaseChanged、onModelSimulationFinished、onModelFoulOccurred），但有些状态变更（如模拟进行中的每帧位置更新）不应触发 View 的完整重绘（会导致闪烁）。通过在 onModelSimulationFinished 中集中刷新，而非逐帧推送属性变更，保证了界面的流畅性。

\section{个人心得与感悟}

Model 层的开发最核心的挑战不在于实现物理公式——这些在课本上都有——而在于处理“不完美的现实”。球-球碰撞的恢复系数、库边的不同弹性、摩擦力的非线性表现，这些细节在理论推导中被简化，但在实际模拟中会通过视觉反馈暴露出来。多次的参数试错让我理解了“物理引擎调优”是一门需要耐心和反复实验的工作。

ViewModel 层的设计让我体会到接口设计的权衡：暴露太少属性，View 层无法获取足够信息；暴露太多，性能和维护成本上升。Q\_PROPERTY + NOTIFY 的组合是 Qt 提供的优雅解决方案，但需要仔细设计每个属性的更新时机以避免不必要的重绘。

与 A 和 C 的跨层协作中，最大的收获是理解了“约定优于实现”——只要层间接口（信号签名、属性名称）提前约定清楚，三方可以真正并行工作。B 在 Model 层增加 foulOccurred 信号时，C 只需在 App 层增加一条 connect 即完成集成，这就是好的接口设计带来的效率。

\section{待完成}

\begin{itemize}
\item 完整斯诺克规则（合法击球判定、所有犯规场景完备、自由球）
\item 彩球按序清空逻辑
\item 输赢判定（红球清空 + 彩球按序 + 分差不可追）
\item Model 层单元测试
\item 边界情况测试（极端角度、零力度、贴库球）
\end{itemize}

\end{document}
